% Authors: Dr. Denys Dutykh (Mathematics Department, Khalifa University of Science and Technology,
% Abu Dhabi, UAE) and Prof. Laurent Vuillon (Univ. Savoie Mont Blanc, CNRS, LAMA, Chambery, France)
\section{A representation problem attached to Markov numbers}

The Markov equation
$$
 x^2+y^2+z^2=3xyz
$$
has a tree of positive integer solutions generated from $(1,1,1)$ by permutations and
the replacement $(x,y,z)\mapsto(x,y,3xy-z)$. The Frobenius--Markov uniqueness
conjecture asks whether the largest entry determines a solution up to
permutation. A useful approach is to attach an arithmetic optimization
problem to each occurrence of that entry in the tree. If two occurrences
are both optimal, a known ordering theorem can distinguish them.

We use the Farey indexing $m_{P/Q}$ of Markov numbers, with
$0<P<Q$ and $\gcd(P,Q)=1$, as in \cite{Aigner2013} and
\cite{LagisquetEtAl2021}. An occurrence includes its reduced fraction;
its label is the integer $m_{P/Q}$. Keeping these notions separate is
essential in a uniqueness problem. The first labels in our convention
are $m_{1/2}=5$, $m_{1/3}=13$, and $m_{2/3}=29$.

Every such occurrence supplies a primitive positive representation
$m_{P/Q}=r^2+s^2$. The subtractive Euclidean algorithm gives this pair
an integer cost $\ell(r,s)$: count single subtractions until $(1,1)$ is
reached. The cost can be read in two equivalent ways, and the paper uses
both. If $s/r=[q_0;\ldots,q_j]$ is the canonical finite continued
fraction, with last quotient at least two, the cost is the sum of the
partial quotients minus one, $q_0+\cdots+q_j-1$. In matrix form,
$$
 \binom sr=P_{q_0}\cdots P_{q_{j-1}}P_{q_j-1}\binom11,\qquad
 P_q=\begin{pmatrix}q&1\\1&0\end{pmatrix},
$$
and the cost is the sum of the indices of the factors;
\cref{lem:cost-descriptions} proves that the two readings agree. The
occurrence pair has cost $P+Q-2$; this follows
from the continued fraction of \c{C}anak\c{c}{\i} and Schiffler
\cite{CanakciSchiffler2020}, as we recall in \cref{sec:occurrences}.
Our question is whether any other primitive representation of the same
integer can have smaller cost. The competitors include all partial
quotients and need not arise from another Markov occurrence.

\begin{theorem}[Thirteen Farey rows]\label{thm:thirteen-rows}
Let $0<P<Q$ be coprime integers with $P\leqslant13$. For every
primitive positive representation
$$
 m_{P/Q}=r^2+s^2,\qquad 0<r<s,
$$
the subtractive Euclidean cost satisfies
$$
 \ell(r,s)\geqslant P+Q-2,
$$
with equality for the occurrence-induced representation.
\end{theorem}

There is no upper bound on $Q$ or on the label. The minimum is weak:
the theorem permits other representations of the same cost. For the
$248$ occurrences of the finite windows of \cref{sec:gaussian}, the
occurrence pair is moreover the only representation of minimum cost.
Since equal labels with equal costs come from the same fraction
(\cref{lem:fixed-cost-injectivity}), weak minimality separates occurrences.

\begin{corollary}\label{cor:thirteen-injective}
If $P/Q$ and $P'/Q'$ are reduced fractions of $(0,1)$ whose numerators are
at most thirteen, then $m_{P/Q}=m_{P'/Q'}$ implies $P/Q=P'/Q'$.
\end{corollary}

\Cref{cor:thirteen-injective} is not new in strength. Lee, Li, Rabideau
and Schiffler show that two distinct fractions with equal labels have a
slope $(P-P')/(Q-Q')$ strictly between $-5/4$ and $-8/7$, and their
numerators then differ by at least six
\cite[Corollary~1.5]{LeeEtAl2023}; this gives the injectivity on the first
six rows, and on the first thirteen rows it leaves sixteen one-parameter
families of candidate pairs, which their fixed-row limits
\cite[Lemma~6.4 and (6.4)]{LeeEtAl2023} (for the first row, the exact value
$m_{1/Q}=F_{2Q+1}$) and a few explicit comparisons of labels exclude. What \cref{cor:thirteen-injective} adds is its
route: it comes from a minimality principle, which makes sense on every row
and would give the uniqueness conjecture if it held on every row. The new
content of the paper is \cref{thm:thirteen-rows} itself, a statement about
every primitive representation of the label, Markov or not, and the method
that proves it.

\subsection{The mechanism of the proof}

Two descriptions of a primitive pair organize the argument. Gaussian
factorization enumerates every representation of one fixed integer.
Inverse Euclidean steps organize every pair of one fixed cost. The
first description treats a finite set of labels completely; the second
controls the unbounded part of each row. \Cref{fig:proof-map} shows how
the two descriptions divide the thirteen rows between them.

At cost $c$, the largest norm is the Fibonacci value $F_{2c+3}$.
Every other path departs from the Fibonacci path at some steps, called
defects. Record the lengths of the Fibonacci runs between successive
defects. The main general tool proved here, \cref{thm:tail-parity}, is
that, on fixing their parities, the full normalized norm $N/F_{2c+3}$ is
strictly monotone in each run length, with an explicit direction. This
holds for every positive number of defects, not only for the strata used
in the thirteen-row theorem. It reduces each infinite box of run lengths to
two exact corner values in $\Q(\sqrt5)$.

An exact expansion of the Cohn trace, in the eigenbasis of the first Cohn
matrix, places the source labels of each row in narrow normalized
intervals around an explicit constant, and we show that every shorter path
misses the interval belonging to its row. The covers used at that step admit unbounded run
lengths: their role is to prove uniform inequalities, not to extrapolate
from a finite search. What remains after them is a finite list of labels,
for which the complete two-square fiber is determined by Gaussian
factorization. The finitely many assertions that these two steps leave to
a computer are collected in one proposition, \cref{prop:finite-verifications};
everything else is proved in the text.

\begin{figure}[tb]
\centering
% Authors: Dr. Denys Dutykh (Mathematics Department, Khalifa University of Science and
% Technology, Abu Dhabi, UAE) and Prof. Laurent Vuillon (Univ. Savoie Mont Blanc, CNRS, LAMA)
% Generated by supplement/code/make_figures.py from exact data; do not edit by hand.
\begin{tikzpicture}[font=\footnotesize, >={Stealth[length=4pt]}]
\draw[teal, line width=1.1pt] (1.049,5.65) -- (1.794,5.65);
\node[ink, right] at (1.831,5.65) {norm envelopes};
\fill[wine] (4.625,5.65) circle (1.1pt);
\node[ink, right] at (4.737,5.65) {complete Gaussian fibers};
\draw[navy, line width=1.7pt] (8.648,5.65) -- (9.393,5.65);
\node[ink, right] at (9.43,5.65) {parity-order covers};
\node[ink, left] at (0.8,0) {$1$};
\draw[teal, line width=1.1pt, ->] (1.049,0) -- (13.267,0);
\node[ink, left] at (0.8,0.4) {$2$};
\draw[teal, line width=1.1pt, ->] (1.123,0.4) -- (13.267,0.4);
\node[ink, left] at (0.8,0.8) {$3$};
\draw[teal, line width=1.1pt, ->] (1.198,0.8) -- (13.267,0.8);
\node[ink, left] at (0.8,1.2) {$4$};
\draw[teal, line width=1.1pt, ->] (1.272,1.2) -- (13.267,1.2);
\node[ink, left] at (0.8,1.6) {$5$};
\draw[teal, line width=1.1pt, ->] (1.347,1.6) -- (13.267,1.6);
\node[ink, left] at (0.8,2) {$6$};
\draw[teal, line width=1.1pt, ->] (1.421,2) -- (13.267,2);
\node[ink, left] at (0.8,2.4) {$7$};
\draw[teal, line width=1.1pt, ->] (1.496,2.4) -- (13.267,2.4);
\node[ink, left] at (0.8,2.8) {$8$};
\draw[mist, line width=0.5pt] (1.571,2.8) -- (3.06,2.8);
\fill[wine] (1.571,2.8) circle (1.1pt);
\fill[wine] (1.72,2.8) circle (1.1pt);
\fill[wine] (1.869,2.8) circle (1.1pt);
\fill[wine] (2.018,2.8) circle (1.1pt);
\fill[wine] (2.167,2.8) circle (1.1pt);
\fill[wine] (2.316,2.8) circle (1.1pt);
\fill[wine] (2.465,2.8) circle (1.1pt);
\fill[wine] (2.614,2.8) circle (1.1pt);
\fill[wine] (2.762,2.8) circle (1.1pt);
\fill[wine] (2.911,2.8) circle (1.1pt);
\draw[navy, line width=1.7pt, ->] (3.06,2.8) -- (13.267,2.8);
\draw[ink, line width=0.6pt] (3.06,2.67) -- (3.06,2.93);
\node[ink, above right, inner sep=1pt, font=\scriptsize] at (3.06,2.82) {$29$};
\node[ink, left] at (0.8,3.2) {$9$};
\draw[mist, line width=0.5pt] (1.645,3.2) -- (3.507,3.2);
\fill[wine] (1.645,3.2) circle (1.1pt);
\fill[wine] (1.72,3.2) circle (1.1pt);
\fill[wine] (1.869,3.2) circle (1.1pt);
\fill[wine] (1.943,3.2) circle (1.1pt);
\fill[wine] (2.092,3.2) circle (1.1pt);
\fill[wine] (2.167,3.2) circle (1.1pt);
\fill[wine] (2.316,3.2) circle (1.1pt);
\fill[wine] (2.39,3.2) circle (1.1pt);
\fill[wine] (2.539,3.2) circle (1.1pt);
\fill[wine] (2.614,3.2) circle (1.1pt);
\fill[wine] (2.762,3.2) circle (1.1pt);
\fill[wine] (2.837,3.2) circle (1.1pt);
\fill[wine] (2.986,3.2) circle (1.1pt);
\fill[wine] (3.06,3.2) circle (1.1pt);
\fill[wine] (3.209,3.2) circle (1.1pt);
\fill[wine] (3.284,3.2) circle (1.1pt);
\fill[wine] (3.433,3.2) circle (1.1pt);
\draw[navy, line width=1.7pt, ->] (3.507,3.2) -- (13.267,3.2);
\draw[ink, line width=0.6pt] (3.507,3.07) -- (3.507,3.33);
\node[ink, above right, inner sep=1pt, font=\scriptsize] at (3.507,3.22) {$35$};
\node[ink, left] at (0.8,3.6) {$10$};
\draw[mist, line width=0.5pt] (1.72,3.6) -- (3.805,3.6);
\fill[wine] (1.72,3.6) circle (1.1pt);
\fill[wine] (1.869,3.6) circle (1.1pt);
\fill[wine] (2.167,3.6) circle (1.1pt);
\fill[wine] (2.316,3.6) circle (1.1pt);
\fill[wine] (2.465,3.6) circle (1.1pt);
\fill[wine] (2.614,3.6) circle (1.1pt);
\fill[wine] (2.911,3.6) circle (1.1pt);
\fill[wine] (3.06,3.6) circle (1.1pt);
\fill[wine] (3.209,3.6) circle (1.1pt);
\fill[wine] (3.358,3.6) circle (1.1pt);
\fill[wine] (3.656,3.6) circle (1.1pt);
\draw[navy, line width=1.7pt, ->] (3.805,3.6) -- (13.267,3.6);
\draw[ink, line width=0.6pt] (3.805,3.47) -- (3.805,3.73);
\node[ink, above right, inner sep=1pt, font=\scriptsize] at (3.805,3.62) {$39$};
\node[ink, left] at (0.8,4) {$11$};
\draw[mist, line width=0.5pt] (1.794,4) -- (6.562,4);
\fill[wine] (1.794,4) circle (1.1pt);
\fill[wine] (1.869,4) circle (1.1pt);
\fill[wine] (1.943,4) circle (1.1pt);
\fill[wine] (2.018,4) circle (1.1pt);
\fill[wine] (2.092,4) circle (1.1pt);
\fill[wine] (2.167,4) circle (1.1pt);
\fill[wine] (2.241,4) circle (1.1pt);
\fill[wine] (2.316,4) circle (1.1pt);
\fill[wine] (2.39,4) circle (1.1pt);
\fill[wine] (2.465,4) circle (1.1pt);
\fill[wine] (2.614,4) circle (1.1pt);
\fill[wine] (2.688,4) circle (1.1pt);
\fill[wine] (2.762,4) circle (1.1pt);
\fill[wine] (2.837,4) circle (1.1pt);
\fill[wine] (2.911,4) circle (1.1pt);
\fill[wine] (2.986,4) circle (1.1pt);
\fill[wine] (3.06,4) circle (1.1pt);
\fill[wine] (3.135,4) circle (1.1pt);
\fill[wine] (3.209,4) circle (1.1pt);
\fill[wine] (3.284,4) circle (1.1pt);
\fill[wine] (3.433,4) circle (1.1pt);
\fill[wine] (3.507,4) circle (1.1pt);
\fill[wine] (3.582,4) circle (1.1pt);
\fill[wine] (3.656,4) circle (1.1pt);
\fill[wine] (3.731,4) circle (1.1pt);
\fill[wine] (3.805,4) circle (1.1pt);
\fill[wine] (3.88,4) circle (1.1pt);
\fill[wine] (3.954,4) circle (1.1pt);
\fill[wine] (4.029,4) circle (1.1pt);
\fill[wine] (4.104,4) circle (1.1pt);
\fill[wine] (4.253,4) circle (1.1pt);
\fill[wine] (4.327,4) circle (1.1pt);
\fill[wine] (4.402,4) circle (1.1pt);
\fill[wine] (4.476,4) circle (1.1pt);
\fill[wine] (4.55,4) circle (1.1pt);
\fill[wine] (4.625,4) circle (1.1pt);
\fill[wine] (4.699,4) circle (1.1pt);
\fill[wine] (4.774,4) circle (1.1pt);
\fill[wine] (4.848,4) circle (1.1pt);
\fill[wine] (4.923,4) circle (1.1pt);
\fill[wine] (5.072,4) circle (1.1pt);
\fill[wine] (5.147,4) circle (1.1pt);
\fill[wine] (5.221,4) circle (1.1pt);
\fill[wine] (5.296,4) circle (1.1pt);
\fill[wine] (5.37,4) circle (1.1pt);
\fill[wine] (5.445,4) circle (1.1pt);
\fill[wine] (5.519,4) circle (1.1pt);
\fill[wine] (5.594,4) circle (1.1pt);
\fill[wine] (5.668,4) circle (1.1pt);
\fill[wine] (5.742,4) circle (1.1pt);
\fill[wine] (5.891,4) circle (1.1pt);
\fill[wine] (5.966,4) circle (1.1pt);
\fill[wine] (6.04,4) circle (1.1pt);
\fill[wine] (6.115,4) circle (1.1pt);
\fill[wine] (6.189,4) circle (1.1pt);
\fill[wine] (6.264,4) circle (1.1pt);
\fill[wine] (6.338,4) circle (1.1pt);
\fill[wine] (6.413,4) circle (1.1pt);
\fill[wine] (6.487,4) circle (1.1pt);
\draw[navy, line width=1.7pt, ->] (6.562,4) -- (13.267,4);
\draw[ink, line width=0.6pt] (6.562,3.87) -- (6.562,4.13);
\node[ink, above right, inner sep=1pt, font=\scriptsize] at (6.562,4.02) {$76$};
\node[ink, left] at (0.8,4.4) {$12$};
\draw[mist, line width=0.5pt] (1.869,4.4) -- (8.425,4.4);
\fill[wine] (1.869,4.4) circle (1.1pt);
\fill[wine] (2.167,4.4) circle (1.1pt);
\fill[wine] (2.316,4.4) circle (1.1pt);
\fill[wine] (2.614,4.4) circle (1.1pt);
\fill[wine] (2.762,4.4) circle (1.1pt);
\fill[wine] (3.06,4.4) circle (1.1pt);
\fill[wine] (3.209,4.4) circle (1.1pt);
\fill[wine] (3.507,4.4) circle (1.1pt);
\fill[wine] (3.656,4.4) circle (1.1pt);
\fill[wine] (3.954,4.4) circle (1.1pt);
\fill[wine] (4.104,4.4) circle (1.1pt);
\fill[wine] (4.402,4.4) circle (1.1pt);
\fill[wine] (4.55,4.4) circle (1.1pt);
\fill[wine] (4.848,4.4) circle (1.1pt);
\fill[wine] (4.998,4.4) circle (1.1pt);
\fill[wine] (5.296,4.4) circle (1.1pt);
\fill[wine] (5.445,4.4) circle (1.1pt);
\fill[wine] (5.742,4.4) circle (1.1pt);
\fill[wine] (5.891,4.4) circle (1.1pt);
\fill[wine] (6.189,4.4) circle (1.1pt);
\fill[wine] (6.338,4.4) circle (1.1pt);
\fill[wine] (6.636,4.4) circle (1.1pt);
\fill[wine] (6.785,4.4) circle (1.1pt);
\fill[wine] (7.083,4.4) circle (1.1pt);
\fill[wine] (7.232,4.4) circle (1.1pt);
\fill[wine] (7.53,4.4) circle (1.1pt);
\fill[wine] (7.679,4.4) circle (1.1pt);
\fill[wine] (7.978,4.4) circle (1.1pt);
\fill[wine] (8.127,4.4) circle (1.1pt);
\draw[navy, line width=1.7pt, ->] (8.425,4.4) -- (13.267,4.4);
\draw[ink, line width=0.6pt] (8.425,4.27) -- (8.425,4.53);
\node[ink, above right, inner sep=1pt, font=\scriptsize] at (8.425,4.42) {$101$};
\node[ink, left] at (0.8,4.8) {$13$};
\draw[mist, line width=0.5pt] (1.943,4.8) -- (11.777,4.8);
\fill[wine] (1.943,4.8) circle (1.1pt);
\fill[wine] (2.018,4.8) circle (1.1pt);
\fill[wine] (2.092,4.8) circle (1.1pt);
\fill[wine] (2.167,4.8) circle (1.1pt);
\fill[wine] (2.241,4.8) circle (1.1pt);
\fill[wine] (2.316,4.8) circle (1.1pt);
\fill[wine] (2.39,4.8) circle (1.1pt);
\fill[wine] (2.465,4.8) circle (1.1pt);
\fill[wine] (2.539,4.8) circle (1.1pt);
\fill[wine] (2.614,4.8) circle (1.1pt);
\fill[wine] (2.688,4.8) circle (1.1pt);
\fill[wine] (2.762,4.8) circle (1.1pt);
\fill[wine] (2.911,4.8) circle (1.1pt);
\fill[wine] (2.986,4.8) circle (1.1pt);
\fill[wine] (3.06,4.8) circle (1.1pt);
\fill[wine] (3.135,4.8) circle (1.1pt);
\fill[wine] (3.209,4.8) circle (1.1pt);
\fill[wine] (3.284,4.8) circle (1.1pt);
\fill[wine] (3.358,4.8) circle (1.1pt);
\fill[wine] (3.433,4.8) circle (1.1pt);
\fill[wine] (3.507,4.8) circle (1.1pt);
\fill[wine] (3.582,4.8) circle (1.1pt);
\fill[wine] (3.656,4.8) circle (1.1pt);
\fill[wine] (3.731,4.8) circle (1.1pt);
\fill[wine] (3.88,4.8) circle (1.1pt);
\fill[wine] (3.954,4.8) circle (1.1pt);
\fill[wine] (4.029,4.8) circle (1.1pt);
\fill[wine] (4.104,4.8) circle (1.1pt);
\fill[wine] (4.178,4.8) circle (1.1pt);
\fill[wine] (4.253,4.8) circle (1.1pt);
\fill[wine] (4.327,4.8) circle (1.1pt);
\fill[wine] (4.402,4.8) circle (1.1pt);
\fill[wine] (4.476,4.8) circle (1.1pt);
\fill[wine] (4.55,4.8) circle (1.1pt);
\fill[wine] (4.625,4.8) circle (1.1pt);
\fill[wine] (4.699,4.8) circle (1.1pt);
\fill[wine] (4.848,4.8) circle (1.1pt);
\fill[wine] (4.923,4.8) circle (1.1pt);
\fill[wine] (4.998,4.8) circle (1.1pt);
\fill[wine] (5.072,4.8) circle (1.1pt);
\fill[wine] (5.147,4.8) circle (1.1pt);
\fill[wine] (5.221,4.8) circle (1.1pt);
\fill[wine] (5.296,4.8) circle (1.1pt);
\fill[wine] (5.37,4.8) circle (1.1pt);
\fill[wine] (5.445,4.8) circle (1.1pt);
\fill[wine] (5.519,4.8) circle (1.1pt);
\fill[wine] (5.594,4.8) circle (1.1pt);
\fill[wine] (5.668,4.8) circle (1.1pt);
\fill[wine] (5.817,4.8) circle (1.1pt);
\fill[wine] (5.891,4.8) circle (1.1pt);
\fill[wine] (5.966,4.8) circle (1.1pt);
\fill[wine] (6.04,4.8) circle (1.1pt);
\fill[wine] (6.115,4.8) circle (1.1pt);
\fill[wine] (6.189,4.8) circle (1.1pt);
\fill[wine] (6.264,4.8) circle (1.1pt);
\fill[wine] (6.338,4.8) circle (1.1pt);
\fill[wine] (6.413,4.8) circle (1.1pt);
\fill[wine] (6.487,4.8) circle (1.1pt);
\fill[wine] (6.562,4.8) circle (1.1pt);
\fill[wine] (6.636,4.8) circle (1.1pt);
\fill[wine] (6.785,4.8) circle (1.1pt);
\fill[wine] (6.86,4.8) circle (1.1pt);
\fill[wine] (6.934,4.8) circle (1.1pt);
\fill[wine] (7.009,4.8) circle (1.1pt);
\fill[wine] (7.083,4.8) circle (1.1pt);
\fill[wine] (7.158,4.8) circle (1.1pt);
\fill[wine] (7.232,4.8) circle (1.1pt);
\fill[wine] (7.307,4.8) circle (1.1pt);
\fill[wine] (7.381,4.8) circle (1.1pt);
\fill[wine] (7.456,4.8) circle (1.1pt);
\fill[wine] (7.53,4.8) circle (1.1pt);
\fill[wine] (7.605,4.8) circle (1.1pt);
\fill[wine] (7.754,4.8) circle (1.1pt);
\fill[wine] (7.829,4.8) circle (1.1pt);
\fill[wine] (7.903,4.8) circle (1.1pt);
\fill[wine] (7.978,4.8) circle (1.1pt);
\fill[wine] (8.052,4.8) circle (1.1pt);
\fill[wine] (8.127,4.8) circle (1.1pt);
\fill[wine] (8.201,4.8) circle (1.1pt);
\fill[wine] (8.275,4.8) circle (1.1pt);
\fill[wine] (8.35,4.8) circle (1.1pt);
\fill[wine] (8.425,4.8) circle (1.1pt);
\fill[wine] (8.499,4.8) circle (1.1pt);
\fill[wine] (8.573,4.8) circle (1.1pt);
\fill[wine] (8.723,4.8) circle (1.1pt);
\fill[wine] (8.797,4.8) circle (1.1pt);
\fill[wine] (8.871,4.8) circle (1.1pt);
\fill[wine] (8.946,4.8) circle (1.1pt);
\fill[wine] (9.021,4.8) circle (1.1pt);
\fill[wine] (9.095,4.8) circle (1.1pt);
\fill[wine] (9.169,4.8) circle (1.1pt);
\fill[wine] (9.244,4.8) circle (1.1pt);
\fill[wine] (9.319,4.8) circle (1.1pt);
\fill[wine] (9.393,4.8) circle (1.1pt);
\fill[wine] (9.467,4.8) circle (1.1pt);
\fill[wine] (9.542,4.8) circle (1.1pt);
\fill[wine] (9.691,4.8) circle (1.1pt);
\fill[wine] (9.765,4.8) circle (1.1pt);
\fill[wine] (9.84,4.8) circle (1.1pt);
\fill[wine] (9.915,4.8) circle (1.1pt);
\fill[wine] (9.989,4.8) circle (1.1pt);
\fill[wine] (10.063,4.8) circle (1.1pt);
\fill[wine] (10.138,4.8) circle (1.1pt);
\fill[wine] (10.213,4.8) circle (1.1pt);
\fill[wine] (10.287,4.8) circle (1.1pt);
\fill[wine] (10.361,4.8) circle (1.1pt);
\fill[wine] (10.436,4.8) circle (1.1pt);
\fill[wine] (10.511,4.8) circle (1.1pt);
\fill[wine] (10.659,4.8) circle (1.1pt);
\fill[wine] (10.734,4.8) circle (1.1pt);
\fill[wine] (10.809,4.8) circle (1.1pt);
\fill[wine] (10.883,4.8) circle (1.1pt);
\fill[wine] (10.957,4.8) circle (1.1pt);
\fill[wine] (11.032,4.8) circle (1.1pt);
\fill[wine] (11.107,4.8) circle (1.1pt);
\fill[wine] (11.181,4.8) circle (1.1pt);
\fill[wine] (11.255,4.8) circle (1.1pt);
\fill[wine] (11.33,4.8) circle (1.1pt);
\fill[wine] (11.405,4.8) circle (1.1pt);
\fill[wine] (11.479,4.8) circle (1.1pt);
\fill[wine] (11.628,4.8) circle (1.1pt);
\fill[wine] (11.703,4.8) circle (1.1pt);
\draw[navy, line width=1.7pt, ->] (11.777,4.8) -- (13.267,4.8);
\draw[ink, line width=0.6pt] (11.777,4.67) -- (11.777,4.93);
\node[ink, above right, inner sep=1pt, font=\scriptsize] at (11.777,4.82) {$146$};
\draw[ink, line width=0.5pt] (0.9,-0.45) -- (13.267,-0.45);
\draw[ink, line width=0.5pt] (0.9,-0.45) -- (0.9,-0.53);
\node[ink, below] at (0.9,-0.53) {$0$};
\draw[ink, line width=0.5pt] (2.762,-0.45) -- (2.762,-0.53);
\node[ink, below] at (2.762,-0.53) {$25$};
\draw[ink, line width=0.5pt] (4.625,-0.45) -- (4.625,-0.53);
\node[ink, below] at (4.625,-0.53) {$50$};
\draw[ink, line width=0.5pt] (6.487,-0.45) -- (6.487,-0.53);
\node[ink, below] at (6.487,-0.53) {$75$};
\draw[ink, line width=0.5pt] (8.35,-0.45) -- (8.35,-0.53);
\node[ink, below] at (8.35,-0.53) {$100$};
\draw[ink, line width=0.5pt] (10.213,-0.45) -- (10.213,-0.53);
\node[ink, below] at (10.213,-0.53) {$125$};
\draw[ink, line width=0.5pt] (12.075,-0.45) -- (12.075,-0.53);
\node[ink, below] at (12.075,-0.53) {$150$};
\node[ink] at (7.083,-1.2) {denominator $Q$};
\node[ink, rotate=90] at (0.15,2.4) {numerator $P$};
\end{tikzpicture}

\caption{How the proof covers the thirteen rows. Each horizontal line is a
Farey row of numerator $P$, drawn against the denominator $Q$. The rows
$P\leqslant7$ are settled by the norm envelopes of \cref{sec:early-rows}
(teal). In the rows $8\leqslant P\leqslant13$ the dots are the $248$ reduced
fractions of the finite windows, settled by complete Gaussian fibers in
\cref{sec:gaussian} (wine); from the junction denominator $Q_P$ on, marked by
a tick, the unbounded tail is settled by the parity-order covers of
\cref{sec:tails} and the comparisons of \cref{sec:assembly} (navy). The arrows
indicate that the rows continue without bound.}
\label{fig:proof-map}
\end{figure}

\Cref{sec:occurrences} gives the two readings of the cost, recalls the
occurrence cost identity and proves the restricted injectivity reduction.
\Cref{sec:gaussian} describes the complete finite fibers and locates the
occurrence pair in them, by one exact Gaussian division per step of the
Farey tree (\cref{prop:vieta-jump}).
\Cref{sec:early-rows} develops the path envelopes and proves the
first seven rows. \Cref{sec:tails} establishes the
general parity order and the construction of the covers.
\Cref{sec:assembly} states the finite verifications, draws their
consequences, joins the domains and proves the theorem. \Cref{sec:outlook}
places the result in the larger question, with three conjectures that imply
the uniqueness conjecture and the computations that support them.

% Authors: Dr. Denys Dutykh (Mathematics Department, Khalifa University of Science and Technology,
% Abu Dhabi, UAE) and Prof. Laurent Vuillon (Univ. Savoie Mont Blanc, CNRS, LAMA, Chambery, France)
\subsection{Relation to earlier work}

The representations studied here belong to an established connection
between Markov numbers, words, and continued fractions. Reutenauer's
Christoffel correspondence parametrizes whole Markov triples
\cite{Reutenauer2009}; the palindromic-closure formula of Reutenauer and
Vuillon gives an occurrence-wise construction of their entries
\cite{ReutenauerVuillon2017}. The underlying central-word and incidence
identities are developed in \cite{BertheDeLucaReutenauer2008}.
These correspondences distinguish occurrences without presupposing that
their largest numerical entries are distinct.

A particularly close arithmetic formulation appears in
\c{C}anak\c{c}{\i} and Schiffler \cite{CanakciSchiffler2020}.
They relate primitive sums of two squares to even palindromic continued
fractions \cite[Theorem~3.10 and Corollary~3.14]{CanakciSchiffler2020}. They
construct the special continued fraction attached to a Markov occurrence
\cite[Theorem~4.1]{CanakciSchiffler2020}, and its right half, through the
proof of \cite[Corollary~4.4]{CanakciSchiffler2020}, is the occurrence pair of
\cref{prop:occurrence-cost} together with its cost $P+Q-2$, as we show after
that proposition.
They conjecture the uniqueness of the representation whose
palindromification is a Markov snake graph
\cite[Conjecture~4.5]{CanakciSchiffler2020} and identify this conjecture
with Markov uniqueness \cite[Theorem~4.8]{CanakciSchiffler2020}. Their
Conjecture~4.6 asks for uniqueness
under the weaker requirement that the half continued fraction use only
the digits $1$ and $2$; they checked both conjectures by computer for all
Markov numbers with denominator below $70$ \cite[Remark~4.7]{CanakciSchiffler2020}.
Our comparison retains every primitive positive representation of the
same integer, with no restriction on its partial quotients. The assertion
is that the occurrence-induced representation minimizes the sum of those
quotients, with the final Euclidean step omitted. It permits ties and
does not assert either of these uniqueness conjectures.

Ordering results compare a different set of objects: the Markov labels
at specified lattice points. Aigner's book records conjectures for rows of
fixed numerator, denominator and sum \cite[Section~10.1]{Aigner2013}; Rabideau and
Schiffler proved the first \cite{RabideauSchiffler2020}, and Lagisquet,
Pelantov\'a, Tavenas, and Vuillon proved all three
\cite{LagisquetEtAl2021}. Their path minimization holds at a fixed lattice
endpoint. Lee, Li, Rabideau, and Schiffler obtain further comparison
inequalities \cite{LeeEtAl2023}, which confine any coincidence of labels to
a narrow cone of slopes; their fixed-row ratio decreases to a constant,
which \cref{sec:row-labels} computes and bounds on both sides from the Cohn
trace, so that each infinite source row lies in a normalized interval. Gaster gives sharp monotonicity ranges using the
stable norm on the modular torus \cite{Gaster2022}. None of these ordering
statements, as stated, minimizes Euclidean cost over the full primitive
two-square fiber of a label. The classical uniqueness theorem for labels
that are prime powers, or twice prime powers, is the case of a fiber with a
single element; \cref{rem:single-representation} derives it from the
results of \cref{sec:occurrences,sec:gaussian}.

Two of our submitted articles treat neighboring questions. The first
studies the same Euclidean cost on the representations produced by a single
Gaussian reflection \cite{DutykhVuillon2026Endpoint}, and contains the cost
identity of \cref{prop:occurrence-cost} in its source dictionary; the second
studies palindromes in the Cohn monoid \cite{DutykhVuillon2026Reflections},
and contains \cref{lem:fixed-cost-injectivity} as the case of equal sums in
the proof of its proposition on the lengths of directing words. Their full texts can be consulted in the repository of
this series, given at the end of the article.

Related extremal problems fix the length and sum of the entries of a
continuant; Kan determines the numerical maximum under these constraints
\cite[Theorem~5]{Kan2012}.\footnote{For sequences of length $t$ and sum $S$,
the published statement assumes $2\leqslant t\leqslant S<nt$, where $n$, a bound
on the entries, belongs to the next theorem of that article; the corrected
version \cite[Theorem~3.5]{Kan2021} states the result for
$2\leqslant t\leqslant S$.}
Our proof derives the required path bounds
directly. In its infinite part, a parity order for the full normalized
ratio reduces whole boxes of gap lengths to exact corner comparisons.
The finitely many remaining labels are settled through their complete
Gaussian factorizations, so that an infinite comparison argument is
joined to a finite piece of exact arithmetic. The resulting thirteen-row
theorem assumes no upper bound on the denominator or on the Markov
value. Its injectivity consequence applies when both occurrences
belong to the proved rows; it does not exclude a collision with an
occurrence outside them.

